\documentclass[10pt,a4paper]{amsart}
\usepackage[margin=19mm]{geometry}
\usepackage{amsmath,amssymb,mathtools}
\usepackage[scr=rsfs,cal=euler]{mathalfa}
\usepackage{xfrac}
\usepackage[unicode,hidelinks,bookmarks=false]{hyperref}
\newcommand{\ie}{i.e.\ }
\newcommand{\cf}{cf.\ }
\newcommand{\resp}{respectively\ }
\newcommand{\Subsection}{\subsection}
\providecommand{\nobreakdash}{\nobreak\hbox{-}}
\setlength{\emergencystretch}{2em}
\multlinegap0pt
\usepackage{bm}
\usepackage{bbm}
\usepackage{dsfont}
\usepackage{xspace}
\usepackage{cancel}
\usepackage{mathtools}
\usepackage{amsthm}
\makeatletter
\@ifundefined{thm}{\newtheorem{thm}{Thm}}{}
\@ifundefined{lem}{\newtheorem{lem}[thm]{Lemma}}{}
\makeatother
\DeclareMathOperator{\Li}{Li}
\title[On the fundamental solutions of two nonlocal parabolic equat]{On the fundamental solutions of two nonlocal parabolic equations related to logarithmic Laplacians}
\author{Bart Rosenzweig, Jonathan Stanfill}
\date{Source: arXiv:2606.04225v1 (CC BY 4.0)}
\begin{document}
\maketitle

\noindent\emph{Source and licence.} On the fundamental solutions of two nonlocal parabolic equations related to logarithmic Laplacians. Version arXiv:2606.04225v1 (CC BY 4.0), distributed under the Creative Commons Attribution 4.0 International licence, as recorded in the arXiv licence metadata for this exact version and shown on the abstract page https://arxiv.org/abs/2606.04225v1. Full rights notice: licenses/arxiv-2606.04225-CC-BY-4.0.txt.\par\smallskip

\noindent\textit{Editorial scope.} Counted unit: the Lem whose source label is \texttt{Lem:Sqrt} in the article (source0.tex lines 638--647, source label \texttt{Lem:Sqrt}), delivered together with the article's own complete original proof of it (source0.tex lines 656--695, one \texttt{proof} environment of 40 lines). No mathematics is written, repaired or re-derived: statement and proof are the article's own text line for line. The proof is detached from the statement in the source: the article states the result at lines 638--647 and prints its proof at lines 656--695, and the proof environment's own header names the statement, so the association is the article's own. Required context retained uncounted: none. Every definition, display and label of those lines is retained unchanged. Omitted from this package and not redistributed: 1--637, 648--655, 696--886. Nothing inside a retained excerpt is omitted or altered: every retained body is delivered exactly as the article's lines are written, so no omission receipt of any kind accompanies this package. Citations: the retained excerpts carry no citation command at all, so no bibliography entry of the article is transcribed and no reference is redistributed. Reference adaptation: none was needed and none was made; every reference inside the delivered unit resolves inside it, so no unresolved reference or citation is produced. Printed slips kept literal: no printed word, symbol, hypothesis or number of the counted statement or of its proof was altered. Layout and class: the article's own class is \texttt{amsart} with packages including inputenc, amsmath, amssymb, amsthm, bbm, geometry, comment, hyperref, mathtools, stmaryrd, enumerate, appendix, graphicx, csquotes; the wrapper is the lane's pinned \texttt{amsart} editorial wrapper at 10pt on A4, which restates the theorem-like environments the retained text needs and adds the article's own macro definitions that the retained text needs (\texttt{\textbackslash Li}), with extra wrapper packages (bm, bbm, dsfont, xspace, cancel, mathtools). The delivered numbering is the unit's own. Assets: no retained span contains a figure environment, a graphics-inclusion command, a PDF-page inclusion, a table or a TikZ picture; no image file is copied into this package, the package declares no asset dependency and no image file is redistributed with it. Licence: version arXiv:2606.04225v1 of this article is distributed under the Creative Commons Attribution 4.0 International licence, as recorded in the arXiv licence metadata for this exact version and shown on the abstract page https://arxiv.org/abs/2606.04225v1. A full rights notice is delivered at licenses/arxiv-2606.04225-CC-BY-4.0.txt. Characters: every retained excerpt is pure ASCII, so the delivered engine, XeLaTeX with its default Latin Modern text font, reports no missing character.
\begin{lem}\label{Lem:Sqrt}
The following holds for $\alpha\in(0,\pi)$ and $s\geq0$: 
\begin{align}\label{Eq:Sqrtseries}
&Q_\alpha(s)=\frac{1}{\sqrt{1-2e^{-s}\cos(2\alpha)+e^{-2s}}}=\sum_{k\in\mathbb{N}_0}c_k(\alpha)s^k,\quad c_0(\alpha)=\frac{1}{2\sin(\alpha)},\\
&c_k(\alpha)=\begin{cases}
\sin(\alpha)\big(-[i\cot(2\alpha)]^{\frac{1+(-1)^k}{2}}\frac{1}{k!}\Li_{-k}\big(e^{2i\alpha}\big)-\sum_{n=1}^{k-1}c_n(\alpha) c_{k-n}(\alpha)\big),& \alpha\in(0,\pi)\backslash\{\tfrac{\pi}{2}\},\\[2mm]
\frac{(2^{k+1}-1)B_{k+1}}{(k+1)!},& \alpha=\tfrac{\pi}{2},
\end{cases}\ k\in\mathbb{N}.\notag
\end{align}
\end{lem}

\begin{proof}[Proof of Lemma \ref{Lem:Sqrt}.]
We begin with the case $\alpha=\tfrac{\pi}{2}$:
\begin{equation}
Q_{\tfrac{\pi}{2}}(s)=\frac{1}{1+e^{-s}}=\frac{1}{2}+\frac{1}{2}\tanh(\tfrac{s}{2})=\frac{1}{2}+\sum_{n\in\mathbb{N}}
\frac{(2^{2n}-1)B_{2n}}{(2n)!}\, s^{2n-1}.
\end{equation}


For the general case $\alpha\in(0,\pi)\backslash\{\tfrac{\pi}{2}\}$, we consider the square of $Q_\alpha(s)$ to write
\begin{align}
\label{eq:2.49}
\frac{2i\sin(2\alpha)}{1-2e^{-s}\cos(2\alpha)+e^{-2s}}&=\frac{e^{2i\alpha}}{1-e^{-s}e^{2i\alpha}}-\frac{e^{-2i\alpha}}{1-e^{-s}e^{-2i\alpha}}=-e^{2i\alpha}\frac{v_-}{1-v_-}+e^{-2i\alpha}\frac{v_+}{1-v_+},
\end{align}
where we have made the substitution $v_\pm=e^{s\pm 2i\alpha}$ noting that $\partial_s^{j}=(v_\pm\partial_{v_\pm} )^{j}$ under this substitution with $s=0$ corresponding to $v_\pm=e^{\pm 2i\alpha}$. Next, recall the identity (noting $\Li_1(z)=-\ln(1-z)$)
\begin{align}\label{Eq:Polylogrec}
\Li_{-n}(z)=\bigg(z\frac{\partial}{\partial z}\bigg)^n \frac{z}{1-z},\quad n\in\mathbb{N}_0,\ z\neq 1.
\end{align}
A short calculation now yields that (noting for $z=e^{2i\alpha}$, $z\partial_z=\frac{1}{2i}\partial_\alpha$ and $z/(1-z)=-\frac{1}{2}+\frac{i}{2}\cot(\alpha)$)
\begin{align}\label{Eq:Polylogid}
\Li_{-j}(e^{2i\alpha})=\frac{i^{-j}}{2^{j}}\frac{\partial^{j}}{\partial\alpha^{j}}\frac{e^{2i\alpha}}{1-e^{2i\alpha}}=i\frac{i^{-j}}{2^{j+1}}\frac{\partial^{j}}{\partial\alpha^{j}}\cot(\alpha),\quad j\in\mathbb{N},\ \alpha\in(0,\pi),
\end{align}
showing $\Li_{-j}\big(e^{2i\alpha}\big)$ is purely imaginary whenever $j\geq1$ is even, and purely real whenever $j\geq1$ is odd.
Thus we can conclude that the coefficients for the square of the function $Q_\alpha(s)$ can be written as
\begin{align}\label{Eq:akmid}
a_k(\alpha)=\frac{i}{2\sin(2\alpha) k!}\big[e^{2i\alpha}\Li_{-k}\big(e^{-2i\alpha}\big)-e^{-2i\alpha}\Li_{-k}\big(e^{2i\alpha}\big)\big]=\frac{1}{\sin(2\alpha) k!}\Im\big[e^{-2i\alpha}\Li_{-k}\big(e^{2i\alpha}\big)\big],
\end{align}
where we have used that for real $s$ and $z$ not on the branch cut $[1,\infty)$ one has that $\overline{\Li_s(z)}=\Li_s(\overline{z})$.
Thus separating by parity in $k$ leads to
\begin{equation}\label{Eq:ak}
a_k(\alpha)=\begin{cases}
-\frac{1}{k!}\Li_{-k}\big(e^{2i\alpha}\big),& k\geq1 \text{ odd},\\
\frac{-i\cot(2\alpha)}{ k!}\Li_{-k}\big(e^{2i\alpha}\big),& k\geq1 \text{ even},
\end{cases}=-[i\cot(2\alpha)]^{\frac{1+(-1)^k}{2}}\frac{1}{k!}\Li_{-k}\big(e^{2i\alpha}\big),\quad k\in\mathbb{N}.
\end{equation}
Returning to our original square root function $Q_\alpha(s)$, we have by the Cauchy product,
\begin{equation}
a_n(\alpha)=\sum_{k=0}^{n}c_k(\alpha) c_{n-k}(\alpha)\Longrightarrow c_0(\alpha)=\frac{1}{2\sin(\alpha)},\quad c_n(\alpha)=\sin(\alpha)\bigg(a_n(\alpha)-\sum_{k=1}^{n-1}c_k(\alpha) c_{n-k}(\alpha)\bigg),\quad n\in\mathbb{N},
\end{equation}
finishing the proof.
\end{proof}

\end{document}
